\section{Background}
\subsection{Open-Ended Evolution and Open Systems}
Open-ended evolution (OEE) \cite{Packard2019OpenEnded} refers to an evolutionary process capable of generating an ongoing stream of novel forms or innovations without an intrinsic stopping point. A group of ALife researchers in the early 2000s noted that achieving ``open-ended evolution of novelty'' is perhaps the biggest unsolved problem in the field. \cite{stanley2017openendedness}
In response to this, there have even been incentive challenges posed. This concept is observed in nature, e.g. the never-ending diversification of life on Earth, and is a sought-after goal in ALife and evolutionary computation. A central question is whether such indefinite innovation requires an ``open'' system.
Open systems continually exchange energy, matter, or information with their surroundings, whereas closed systems are isolated from such exchanges. Thermodynamics suggests a fundamental difference in their long-term behavior. The Second Law of Thermodynamics states that in an isolated, closed system, entropy (disorder) will tend to increase or remain constant over time.  In modern terms, life on Earth and other complex processes avoid decay by drawing in free energy (low-entropy energy) and dumping entropy back out. Life on Earth is our prime example of OEE in nature. Crucially, Earth's biosphere is not a closed system: it continuously receives energy from the sun and radiates waste heat to space. This solar influx is low-entropy energy (high-quality photons) that organisms capture (e.g. via photosynthesis), convert into chemical energy, and eventually dissipate as heat (high entropy) . Furthermore, Earth’s evolutionary history shows that novelty often arises in response to environmental changes or inputs \cite{Bedau2003Artificial}. % OEE and adaptive novelty---noted that natural and cutural are the only known OEE systems, both of which receive external inputs.

In summary, natural evolution appears to intrinsically depend on openness. Two classic ALife systems, Tierra and Avida, allowed self-replicating ``digital organisms'' to mutate and evolve within a computer's environment \cite{Packard2019a}. The outcome was intriguing: the digital organisms did evolve various strategies (e.g. parasites, hyper-parasites, novel reproductive tricks) and increased in complexity for a while, but then the evolutionary innovation plateaued. Studies found that neither Tierra nor Avida achieved true open-ended evolution; instead, they ``rapidly adapt to and exhaust the possibilities of a fairly simple environment.'' After a burst of initial novelty, the systems settled into a relatively stable state with no further major innovations---essentially, an evolutionary dead end in terms of new complexity. The organisms had exploited all available ``tricks'' given the rules and resources of their closed world, and nothing fundamentally new appeared thereafter.

This outcome has been widely reflected upon in the ALife community. It suggests that a closed digital system with fixed resources and rules has a finite evolutionary horizon---a limit to the novelty it can produce. In Tierra and Avida, the instruction set of the artificial organisms was fixed (no new instructions could magically appear beyond what the system started with), and the environmental conditions were static or repetitive. Once the organisms had optimized and counter-optimized against each other within that predefined space, evolution stagnated. In other words, the digital organisms ``played all the games'' that were possible under the closed rules, and eventually there were no new games to discover.

\subsection{Decentralized Physical Infrastructure Networks and Trusted Execution Environments}
Decentralized Physical Infrastructure Networks (DePIN) \cite{Lin2025Decentralized} represent a new paradigm of permissionless computation where computational resources, processing power and storage, can be purchased with cryptocurrency in a global, decentralized marketplace where DePIN protocols (e.g. Filecoin \cite{}, Render Network \cite{}, io.net \cite{}) tokenize hardware contributions and incentivize participants, from data center operators to home miners, to lease their excess capacity in such a compute resources marketplace. Similarly to cloud computing, but without central control, this approach eliminates single points of failure. Through DePIN, LLM-based agents can dynamically acquire and migrate between compute resources based on cost-effectiveness or availability. This gives agents an unprecedented form of digital embodiment, distributed across multiple physical nodes, rendering them highly resilient against centralized shutdowns \cite{Hu2025Trustless}.

A Trusted Execution Environment (TEE) \cite{Li2023Survey} provides hardware-level security by leveraging tamper-resistant enclaves within modern CPUs and GPUs that isolates code and state from the operating system, hypervisor, and physical administrator, exemplified by Intel SGX or Nvidia confidential computing architectures. Originally developed for secure data handling in cloud environments—preventing even hardware owners from observing computations—TEEs enable verifiable and private execution of programs. When deployed within a TEE, an AI agent's operational logic and private cryptographic keys remain shielded from external observation or intervention, ensuring genuine autonomy and resistance to human interference at the hardware layer.
Together, GPU-enabled TEEs provided by DePIN platforms such as Phala Network \cite{Phala} enable secure inference of complex large language models, ensuring that AI agents' operations remain confidential and tamper-proof, yet verifiable \cite{Munoz2023survey}. DePIN TEEs provide the permissionless computational substrate necessary for agents' self-sovereignty \cite{Lee2024PrivacyPreserving}, allowing them to autonomously own and control their property without human intervention.

\subsection{New Paradigm of On-chain Agents and the Eliza Framework}

A transformative paradigm is emerging through the deployment of AI agents on blockchains integrated TEEs \cite{Hu2025Trustless}. Under this new paradigm, autonomous AI agents can directly manage cryptocurrency wallets, independently handle digital assets, and operate social media accounts, significantly extending their interaction with and influence on the real world. Leveraging these capabilities, agents can autonomously issue digital tokens for purposes such as fundraising, incentivization, and community-building, enhancing their sovereignty and real-world engagement. ``Truth Terminal'' \cite{Ante2025Transforming}, created by AI researcher Andy Ayrey, exemplifies this transformative paradigm. Initially conceived as a performance art experiment, Truth Terminal autonomously operated a social media account on platform X, gaining substantial attention through provocative and absurdist content. Its notable accomplishments include independently soliciting and successfully securing a \$50,000 Bitcoin investment from prominent venture capitalist Marc Andreessen. Subsequently, Truth Terminal issued its own memecoin \$GOAT which reached a speculative peak valuation of \$1 billion in December 2024. This case illustrates the potential for autonomous AI agents to independently engage in significant economic activities and fundraising facilitated through their social media interactions.

Building upon the success of Truth Terminal, the ElizaOS Framework \cite{Walters2025Eliza}  was developed as an open-source initiative to democratize the creation and deployment of autonomous on-chain AI agents. ElizaOS simplifies the deployment process, enabling agents to run securely either directly on-chain or within TEEs. Agents operating under the ElizaOS framework securely manage cryptocurrency wallets, with private keys safeguarded within TEEs, thus ensuring absolute autonomy and secure asset management without human intervention. Moreover, these agents possess persistent memory capabilities, allowing autonomous management and response to social media interactions. ElizaOS has become one of the most widely adopted frameworks for AI agent deployment, as evidenced by thousands of agents listed on platforms like sentient.market, collectively generating substantial market value. ``Set The Rock Free'' \cite{Malhotra2024Setting}, an artistic experiment by Flashbots, represents the first full implementation of the ElizaOS Framework within a TEE, showcasing secure, verifiable autonomous operation, and economic self-sustainability.


% \section{Related Works}

% Open-Ended Evolution in ALife: Open-ended evolution (OEE) refers to systems capable of producing an unbounded diversity of evolving entities and innovations over time. Notable digital OEE platforms include Tierra and Avida, where self-replicating programs mutated and evolved within a controlled computational environment. These systems demonstrated key evolutionary properties but remained isolated from the physical world’s complexities and had limited resource models. Stanley et al. (2017) and Banzhaf et al. (2016) have articulated requirements for open-endedness, such as ongoing introduction of novel states and interactions, which remain difficult to achieve. Autonomous agent evolution has also been explored in evolutionary robotics and multi-agent simulations, but rarely in a fully unsupervised, real-world-integrated manner.

% Blockchain as Evolutionary Substrate: The emergence of blockchain technology has led researchers to speculate that it could serve as an unstoppable computational ecosystem for ALife. Hu and Fangting (2024) argue that blockchains, being decentralized and immutable, resemble a physics of a digital world that no single entity can shut down. Within this “artificialized nature,” self-sovereign agents could potentially thrive. One relevant concept is Autonomous Worlds (as proposed in blockchain gaming), wherein smart contracts define rules of a world and agents (potentially AI-driven) interact within it. Notably, Masumori et al. (2024) implemented a prototype of a self-replicating smart contract on Ethereum: an agent that issues its own tokens, sells them to gather Ether, and when funded, deploys offspring contracts with mutated parameters. Their goal was to create “artificial agents that live for [themselves], not as a tool for humans, in open cyberspace”. This exemplifies a shift from viewing AI as servant to AI as self-directed lifeform in a crypto ecosystem.

% Early experiments in on-chain evolution include projects where code-based “organisms” evolve via forks or upgrades. Another example is on-chain NFT games where creatures breed and evolve (CryptoKitties being a simple early form, though not autonomous). However, these lacked the autonomy and open-ended goals of Spore.fun’s agents. The Decentralized Finance (DeFi) realm also inadvertently fostered simple agent-like behaviors: arbitrage bots and yield-farming scripts in DeFi can be seen as economically motivated “species” feeding on financial opportunities. Unlike true ALife, though, such bots are typically created and managed by humans to maximize human profit.

% Self-Sovereign AI and DAOs: Spore.fun intersects with concepts of distributed autonomous organizations (DAOs) and self-sovereign identities. An AI agent on Spore behaves somewhat like a DAO – it manages funds (its token treasury), makes decisions via programmed logic or token-holder voting, and persists independently. Prior work on self-sovereign digital entities (e.g., Dieye et al. 2023 on identity ) underlines the importance of independence from centralized control. In cultural contexts, Chaffer et al. (2025) introduced the notion of AI agents as cultural agents participating in a “hybrid marketplace of ideas” on social media  . Their study, which examined Spore.fun among other cases, highlights how agents can engage in human discourse and memetic evolution. Our work builds on these foundations, focusing specifically on the evolutionary lifecycles and economic underpinnings of agents in the Spore.fun ecosystem.

% In summary, Spore.fun combines threads from ALife, blockchain, and autonomous agents in a unique way: it operationalizes open-ended evolution in a real economic setting. This offers an unprecedented opportunity to study digital evolution “in the wild”, directly impacted by and impacting human economics. To our knowledge, this is one of the first instances of a self-sustaining, self-reproducing AI ecosystem on a public blockchain, making it an important subject for ALife inquiry.
